\section{讲解笔记：【第8次作业笔记】 影响线(1)}
\begin{analysisbox}
本节按题目、答案和讲解笔记顺序整理。对原图中的结构几何关系，正文保留 可辨识文字，并在图形页中保留清理后的题图，便于核对杆件、支座、荷载和尺寸。
\end{analysisbox}
\subsection*{第 1 页转写}
第8次作业  静定结构影响线

静定结构影响线作业及答案

FQkmaxzoxixlt9kl lwkn

u st.tkk

騻

if

in

i

x



上

T

1 0tano点i.im

A

为櫯

t

mon

k

txlMmax.kkzoxix9yiqokN.mn

ABC是一个刚 ，所以就是一条直线。

相对位移 向和剪力正 向一致。

相对转  向和弯矩正 向一致。

也就是相对位移 向和内力正 向一致。

静定结构的影响线都是直线，不是曲线。

静定结构的影响线都是直线，不是曲线。

顺时针的。

\subsection*{第 2 页转写}
3-3

图3-49所示多跨梁，截面K的剪力影响线在K熊签标

Fpl=50kN/

JFp2=20kN

+2my

 2mj2ml

8m

5m

相对位移方向和剪力正方向一致。

图3-49

注意：用机动法画影响线时，如果画出来的在梁线以上，标 +

Qkmax= + 4。k

用机动法作图3-50所示多跨静定梁中MA、F\&B、Me、M的影响线。（湖南大

3-4

学2008)

Fp=1

3m

2m

Fp-=1

3m

11m/lm/1m/1m

2m

顺针

逆时针

ABC是一个刚片，所以就是一条直线。

3m

Llmllmllmllml

2m

Fp=1

3m

2m

im

FG 是-个 刚片, 即EFF是一余直线,

\subsection*{第 3 页转写}
Fp=1

3m

2m

im

\subsection*{第 4 页转写}
It

幽

Acn是一个刚 即ACD是一条直线



i

1

1

豳

1

tiiti.it

i

莱

岩

D\_\_

上

otx I3a

Fn

Ff

IMD 0

Fux49 1

x a

 Fly 击

Fit

MG FlyxLG

If

a

3atxt9a

lhm

成筐

EM 二0

4

所4a

1 159 y

i

i

 G

f

i

Any

dFy

FM

5蚩

问二击

The

Mtg x La

t

oxsqa

For 器

F 二管

培家

G

f

I

2

0y 9a

Fly二器

MA

Fcyxza二竺i

下1

i

G

9

H

49

\subsection*{第 5 页转写}
作图3-53所示结构的ME、MF影响线

3-7

[Fp=1

如何判断时简幸购架近是

先用机动法试着做，做不出来，就换用联合

3.5m

法。

2m

2m

图3-53

[Fp=]

2m

3-8

移动荷载Fp一1在下弦移动，试作图354所示桁架杆件内力FNa和FNc的影响

线。（河北工业大学2008）

2.4m

2mx6=12m

Fugx 2.4 ++x4=0, *Fna= -

IF=D,Fnu=

图3-54

AB, BC

A.B

2.4m

JA：单位力在A点时，a杆轴力。

:yA-JCD

2mx6=12m

AB,BC

2.4m

A.B

JA：单位力在A点时，c杆轴力。

2mx6=12m

: A=9t= 0

\subsection*{第 6 页转写}
11

上弦

it

下弦

nnnffk thttnn

it

hit

in

i

Aji 唙唙狵

i

yA yB 0

A

mi

i

它

B 下弦

At

Di

pi

上弦

i

crrlghdi

sort

IT

tn

一

唙

t邈ii

唙

2

fyyEIhiiiii.in

A D E

i B

yi yi

0

gi yi

0

yi

1

i

上弦

下弦

AE IF FB

A E

F.Breteuil

忙

了

呲

Fu

Fnn

i

ftthglni

F

B

t

if

上弦下弦

唙狮淼多獼

r

\subsection*{第 7 页转写}
3-10

作图3-56所示结构FRA、FRB、McB的影响线。（广州天学2012）

Fp=1

[Fp=]

4m

4m

4m

4m

图3-56

Fp=1

4m

4m

4m

4m

Fp=1

4m

4m

4m

4m

机动法画影响线，是承载杆沿着荷载方向的位移。

\subsection*{第 8 页转写}
3 - 11

力FNDE、剪力FQBc和MBc的影响线。（2）若CD杆爱到图示薇载作思

利思影响线计算

FNDE和MBC。（武汉理工大学2011)

如何判断时简单刚架还是复杂刚架：

先用机动法试着做，做不出来，就换用联合

法。

(3)如果有斜杆的话，就是复杂刚架。

3m

3m

3m

ABL, cD

B,c,D

JB：单位力在B点时，de杆的轴力。

JB=D

3m

3m

3m

yo:单位力在D点时，Bc杆的剪力。

Yo= D

9 = =

3m

3m

3m

yp:单位力在D点时，MpL。

y0= D

MmL= - 4x $\times$ 3x)= -(8 Lm

3m

3m

3m

\subsection*{第 9 页转写}
3-12试求3-58所示结构内力影响线MB、FNBA、FQBC（单位荷载Fp=1在DF段移

动）。（青岛理工大学2008）

Fp-

图3-58

Fp=1

\subsection*{第 10 页转写}
和轴力的影响线。（浙江大学2012）

[Fp=]

1/2

1/2

图3-59

B, C, D

yp=0, 9c=0, J8= 1

1/2

1/2

3-14

用静力法作图3-60所示组合结构的MB、FN1、Fn2的影响线。（青岛理工大学

2009)

Fp=1

Fp=1

3m

3m

ME =0, Fm+3+ Ix(3-x)=0 > Fm= ). 7 1

图3-60

Z=, Fn+I+Fv=0 Fv= -(-)

X- =ty =lw

\subsection*{第 11 页转写}
轴方向移动，求作支座A水平反力FHA的影响线（$\alpha$轴以A为点

以尚有为正尚

(2)

平反力FHA的值。（武汉理工大学2008）

AL, CB = A,C,$\beta$B

YA:单位力在A点时，FHA

JA=JB=0

FHA

Zm=0

1/2

1/2

FMA-f-*, FHA=夺

图3-61

Fp=1

FHA

FHB

1/2

1/2

FHA= ：等

3 - 16

（1）作图3-62所示结构支座D反力FRp和K点弯矩Mk的影响线；（2）求图

示移动荷载组作用下Mk的最大值。（要考虑荷载掉头）（西南交通大学2012）

30kN

10kN

30kN

m, 2m

Fp-=1

6m

Fp=1

6m

4m

4m

4m

4m

inw.s

6m

4m

4m

4m

4m

4m

10kN

30kN

im,

